\section{International processing split}\label{sec:upsplit}

The construction-era operations tables assume IN2P3 50\% and UKDF 25\% of DRP
processing, flat across the survey. The current working shares are
\emph{provisional}:

\begin{itemize}
\item \emph{France (CC-IN2P3)}: \smFrShare{} throughout. The facility's own
      capacity plan has not yet been received. CC-IN2P3 has run DP2 pilot
      campaigns, so multi-site DRP has been exercised in practice.
\item \emph{UK (UKDF)}: \smUkShare{} as a work-weighted ten-year aggregate.
      UKDF bought early to cover about 40\% of the first years' requirement,
      and will commit less later so the aggregate is met. The model's
      year-by-year profile meets the aggregate, but it is a placeholder until
      UKDF supplies its own. UKDF's provisioned capacity is \smUkCapM{} million
      core-hours a year.
\item \emph{USDF} carries the residual share of DRP, plus all Alert
      Production, Data Access Centre, Qserv and Kubernetes services.
\end{itemize}

\smpara{Capacity against the current compute basis.}
On the \vthirty{} basis (\secref{sec:upcompute}), UK demand is
\smUkReqDRfour{} of UKDF's provisioned capacity in year 5, \smUkReqDRfive{} in
year 6 and \smUkReqDRnine{} by year 10. The USDF share reaches
\smUsdfSplitDRfive{} of its current fleet in year 6, \smUsdfSplitDRsix{} in year 7
and \smUsdfSplitDRnine{} by year 10
(\tabref{tab:smSplit}). On the superseded \wfortyone{} basis, the UK profile
stayed within its provisioned capacity through year 10. Partner provisioning
should be reviewed against the current estimate.

\smpara{Storage.}
Splitting processing moves live intermediates to the facility doing the work.
Retained releases continue to be archived at the USDF, so the split does not
change the USDF archive.

% GENERATED by generate_model.py from lsst/rubin-sizing-model @ v2026.09.2. Do not edit by hand.
% Regenerate: git clone --branch v2026.09.2 https://github.com/lsst/rubin-sizing-model ; then, inside the clone, uv run generate_model.py --emit-tex <this directory>

\small \begin{longtable}{|p{0.55\textwidth}|r|l|}
\caption{International DRP processing split (provisional). \label{tab:smSplitInputs}}\\
\hline
\textbf{Metric} & \textbf{Value} & \textbf{Unit} \\
\hline\endfirsthead
\hline
\textbf{Metric} & \textbf{Value} & \textbf{Unit} \\
\hline\endhead
France / CC-IN2P3 share & 40\% & \% \\ \hline
UK / UKDF share (10-yr aggregate) & 25\% & \% \\ \hline
USDF share (residual, aggregate) & 35\% & \% \\ \hline
UK provisioned capacity & 40,000,000 & ch/yr \\ \hline
\end{longtable} \normalsize


\begin{landscape}
% GENERATED by generate_model.py from lsst/rubin-sizing-model @ v2026.09.2. Do not edit by hand.
% Regenerate: git clone --branch v2026.09.2 https://github.com/lsst/rubin-sizing-model ; then, inside the clone, uv run generate_model.py --emit-tex <this directory>

\tiny \begin{longtable}{|p{0.22\textwidth}|l|r|r|r|r|r|r|r|r|r|r|}
\caption{DRP work split across the USDF, France (CC-IN2P3) and the UK (UKDF), by loop year. \label{tab:smSplit}}\\
\hline
\textbf{Metric} & \textbf{Unit} & \textbf{DP2 (actual)} & \textbf{DR1} & \textbf{DR2} & \textbf{DR3} & \textbf{DR4} & \textbf{DR5} & \textbf{DR6} & \textbf{DR7} & \textbf{DR8} & \textbf{DR9} \\
 &  & FY2026 & FY2027 & FY2028 & FY2029 & FY2030 & FY2031 & FY2032 & FY2033 & FY2034 & FY2035 \\
\hline\endfirsthead
\hline
\textbf{Metric} & \textbf{Unit} & \textbf{DP2 (actual)} & \textbf{DR1} & \textbf{DR2} & \textbf{DR3} & \textbf{DR4} & \textbf{DR5} & \textbf{DR6} & \textbf{DR7} & \textbf{DR8} & \textbf{DR9} \\
 &  & FY2026 & FY2027 & FY2028 & FY2029 & FY2030 & FY2031 & FY2032 & FY2033 & FY2034 & FY2035 \\
\hline\endhead
USDF share (residual) & \% & 100\% & 20\% & 20\% & 36\% & 36\% & 36\% & 36\% & 36\% & 36\% & 36\% \\ \hline
France / CC-IN2P3 share & \% & 0\% & 40\% & 40\% & 40\% & 40\% & 40\% & 40\% & 40\% & 40\% & 40\% \\ \hline
UK / UKDF share & \% & 0\% & 40\% & 40\% & 24\% & 24\% & 24\% & 24\% & 24\% & 24\% & 24\% \\ \hline
Total DRP core-hours & core-hrs & 2,793,600 & 35,481,751 & 70,963,503 & 106,445,254 & 141,927,006 & 177,408,757 & 212,890,508 & 248,372,260 & 283,854,011 & 319,335,762 \\ \hline
USDF & core-hrs & 2,793,600 & 7,096,350 & 14,192,701 & 38,320,291 & 51,093,722 & 63,867,152 & 76,640,583 & 89,414,013 & 102,187,444 & 114,960,874 \\ \hline
France / CC-IN2P3 & core-hrs & 0 & 14,192,701 & 28,385,401 & 42,578,102 & 56,770,802 & 70,963,503 & 85,156,203 & 99,348,904 & 113,541,604 & 127,734,305 \\ \hline
UK / UKDF & core-hrs & 0 & 14,192,701 & 28,385,401 & 25,546,861 & 34,062,481 & 42,578,102 & 51,093,722 & 59,609,342 & 68,124,963 & 76,640,583 \\ \hline
USDF & cores & 794 & 2,016 & 4,032 & 10,885 & 14,514 & 18,142 & 21,771 & 25,399 & 29,028 & 32,656 \\ \hline
France / CC-IN2P3 & cores & 0 & 4,032 & 8,063 & 12,095 & 16,126 & 20,158 & 24,190 & 28,221 & 32,253 & 36,284 \\ \hline
UK / UKDF & cores & 0 & 4,032 & 8,063 & 7,257 & 9,676 & 12,095 & 14,514 & 16,933 & 19,352 & 21,771 \\ \hline
USDF peak cores vs fleet (Milano-equivalent) & \% & 4\% & 9\% & 19\% & 50\% & 67\% & 84\% & 101\% & 118\% & 134\% & 151\% \\ \hline
UK required vs UK provisioned capacity & \% & 0\% & 35\% & 71\% & 64\% & 85\% & 106\% & 128\% & 149\% & 170\% & 192\% \\ \hline
\end{longtable} \normalsize

\end{landscape}
